\chapter{Discussion}

\label{ch:Discussion}

This chapter builds on the numerical results from Chapter 4 and explores three aspects in depth:
 (i) the physical origins of design performance and its dependence on topological charge;
 (ii) comparison and positioning relative to existing literature;
 (iii) the limitations of this work's methodology and directions for future research.

%===================================================================

\section{Physics of Topology-Dependent Performance}

\label{sec:physics_discussion}

\subsection{Mechanism of SPP-like Geometry Emergence in $\ell=2$ and Free-Form Optimization in Other Modes}

\label{ssec:spiral_emergence}

Chapter 4 showed that, for the $\ell = 2$ design, topology optimization guided by adjoint gradients spontaneously discovers a clear two-arm spiral structure at the center (Figure~\ref{fig:structures_all}(c)). This outcome has a clear physical origin: the azimuthal phase dependence $e^{i2\phi}$ of the LG$_{2,0}$ mode requires the structure to provide a continuously increasing $2 \times 2\pi$ phase difference across $\phi \in [0, 2\pi)$. This requirement is directly met by a two-arm spiral material distribution, which closely mirrors the working principle of a spiral phase plate (SPP). The surrounding scattering features act as a compensation mechanism by which the optimizer corrects near-field coupling errors, further improving the accuracy of far-field mode shaping.

For the other designs ($\ell = 0, 1, 3, 4, 5$), topology optimization does not produce simple $\ell$-fold spiral patterns. Instead, the optimizer converges to free-form, non-intuitive material distributions with no obvious geometric regularity. This reflects the nature of global topology optimization: an SPP-like spiral strategy is not always the globally optimal solution. When the spiral approach does not maximize the figure of merit, the optimizer automatically explores other structural forms. In all designs, the surrounding scattering features become denser as topological charge increases, reflecting the higher modulation demand on azimuthal phase space for higher-order OAM modes.


\subsection{Origin of the Purity--Overlap Decoupling}

\label{ssec:purity_overlap}

Table~\ref{tab:results_summary} reveals a phenomenon worth examining in depth: the Purity for $\ell = 1$ is the highest of all six ($98.66\%$), yet $|{\rm Overlap}|^2$ has already dropped from $60.92$ at $\ell=0$ to $53.88$; for $\ell = 2$, Purity is still $97.96\%$, but $|{\rm Overlap}|^2$ drops further to $7.99$. This "high purity, low overlap" phenomenon arises from the fundamentally different physical meanings of the two metrics:

Purity measures the normalized shape agreement between the far-field pattern and the target LG mode (see equation (3.9)) and is insensitive to the absolute energy level; $|{\rm Overlap}|^2$ quantifies the absolute coupling power, directly reflecting how much of the source radiation is effectively converted into the target mode. As topological charge increases, the ring peak radius of the LG$_{\ell,0}$ mode, $r_{\rm peak} = w_0\sqrt{\ell/2}$, grows rapidly, while the design region size is fixed at $4.5\,\mu{\rm m}$. Effective pattern shaping for high-$\ell$ modes can only happen within this constrained space. The optimizer can produce the correct spiral phase distribution within the finite space (maintaining high Purity), but cannot simultaneously ensure sufficient absolute energy coupling into the larger-scale modes (hence $|{\rm Overlap}|^2$ drops significantly).

This decoupling confirms that using both Purity and $|{\rm Overlap}|^2$ as dual evaluation metrics is necessary: the two are complementary—the former reflects the fidelity of the mode shape, the latter reflects the efficiency of energy conversion, and neither alone is sufficient.

\subsection{Physical Limits of High-Order OAM Generation}

\label{ssec:high_order_limit}

The mode purity for $\ell = 4, 5$ drops to around $82\%$, and slight non-uniformity appears in the far-field phase distribution (see Section~\ref{sec:farfield}). This performance drop can be attributed to two mutually coupled physical limitations:

\textbf{(i) Azimuthal phase gradient resolution limit:} High-order OAM modes require finer azimuthal phase gradients across the design region, with an equivalent azimuthal spatial frequency proportional to $\ell$. At a fixed $4.5\,\mu{\rm m}$ design area and $50\,{\rm grid/}\mu{\rm m}$ resolution, the azimuthal phase modulation needed for high topological charges gradually approaches the spatial frequency ceiling resolvable by the grid, creating a numerical bandwidth limitation—the optimizer can no longer accurately express the required fine geometric phase distribution.

\textbf{(ii) Insufficient design domain size relative to mode scale:} The ring peak radius of LG modes grows proportionally with topological charge, and for $\ell = 4, 5$ the theoretical spot size is approaching the design domain boundary. The optimizer cannot fully express the complete radiation pattern required for high-order modes within the finite design space, limiting far-field mode shaping.

These two limitations together form the performance bottleneck of this work at high $\ell$, and suggest that expanding the design region or increasing numerical resolution would likely significantly improve mode purity for high-order modes. Notably, despite these spatial constraints, the phase winding verification for all designs (Section~\ref{sec:winding}) shows topological charge accuracy errors below $1\%$, indicating that the performance loss at high $\ell$ is primarily expressed as deviations in the intensity profile, not as misalignment of the topological charge itself.

  

%===================================================================

\section{Comparison with Prior Work}

\label{sec:comparison}

\subsection{Metasurface-Based OAM Emitters}

\label{ssec:compare_metasurface}

\begin{table}[htb]
\centering
\caption{Quantitative comparison of this work with representative existing OAM metasurface approaches.}
\label{tab:comparison}
\resizebox{\textwidth}{!}{%
%\begin{tabular}{p{2.8cm}cccp{2.5cm}c}
\begin{tabular}{cccccc}
\hline
Approach & Design method & Material & Purity (\%) & Source type & Valley selectivity \\
\hline
Karimi 2014 \cite{karimi2014generating} & Forward (geometric phase) & Au plasmonic & $\sim$64 & External plane wave & None \\
Guo 2016 \cite{guo2016merging} & Forward (geometric+propagation phase) & Au plasmonic & $>$90 & External plane wave & None \\
Sroor 2020 \cite{sroor2020high} & Forward ($\text{TiO}_2$ metasurface) & $\text{TiO}_2$ all-dielectric & $\sim$90  ($\ell \le 100$) & Intracavity laser & None \\
Li 2023 \cite{li2023valley} & TMD photonic crystal (BIC) & $\text{WSe}_2$/Si & Purity not quantified & \textbf{Active ($\text{WSe}_2$)} & Fixed $\ell$ \\
\textbf{This work} & \textbf{Inverse design (topology optimization)} & \textbf{GaP all-dielectric} & \textbf{82$\sim$ 98 ($\ell \le 3$)} & \textbf{Active ($\text{WS}_2$)} & \textbf{Fixed $\ell$ \& RCP purity $<1\%$} \\
\hline
\end{tabular}
}
\end{table}


Table~\ref{tab:comparison} summarizes a quantitative comparison of this work with representative literature on core metrics. As can be seen, this work holds clear advantages over existing approaches in three dimensions:

\textbf{(i) Simultaneous high purity and active chiral source integration:} Existing approaches either achieve high purity (Sroor 2020: $\sim 90\%$) but use an intracavity metasurface laser, which is a large laser setup rather than emitter integration, and lacks valley selectivity; or combine TMD material  with metasurfaces to act as an active circularly polarized light source (Li 2023) but do not quantify LG mode purity and are limited by fixed lattice geometry. This work achieves LG mode purity ranging from $98.66\%$ ($\ell = 1$) to $94.13\%$ ($\ell = 3$), while simultaneously integrating an active valley-polarized exciton source, filling the "active source + high-purity LG mode" design gap.

\textbf{(ii) Quantitative verification of valley polarization selectivity:} Only a small number of existing works have quantified inter-valley crosstalk. This work directly provides a clear, experimentally benchmarkable baseline through the RCP cross-excitation test ($\ell = 3$ structure: RCP purity $0.03\%$ vs. LCP purity $94.13\%$).

\textbf{(iii) Design freedom:} Using $5.1\times10^4$ design degrees of freedom instead of the traditional few geometric parameters allows the optimizer to automatically search for non-intuitive high-performance structures without any preset geometric shape.

\textbf{(iv) Compact footprint:} The device footprint of this work is $4.5 \times 4.5\;\mu\text{m}^2$. This small size is close to the exciton diffusion length in monolayer $\text{WS}_2$, which helps to reduce intervalley scattering during emission. Topology optimization makes this compact size possible by encoding the full phase profile within a single non-periodic design region. This allows the device to achieve high OAM purity and active valley-locked emission within an ultra-compact footprint.

\subsection{Topology Optimization in Nanophotonics}

\label{ssec:compare_topo}

There are many precedents for applying topology optimization to nanophotonics, including photonic crystal cavity mode shaping and nonlinear frequency conversion efficiency optimization \cite{Hammond2022high}. However, inverse design approaches directly targeting chiral active light sources combined with OAM mode integration remain rare in the existing literature.

The main methodological contributions of this work lie in three points. First, a chiral emitting point source is used as the optimization driving field, so the resulting structure corresponds from the start to the actual radiation environment of valley-polarized excitons. Second, the joint logarithmic objective function of LG mode overlap integral and mode purity (equation (3.16)) is used as the FOM, preventing low-power pathological convergence while maintaining precise control over mode shape. Third, the circular polarization selectivity of the design is maintained throughout, and quantitatively verified through near-field chirality density analysis and polarization cross-excitation testing (Sections~\ref{sec:chirality} and \ref{sec:pol_selectivity}).

%===================================================================

\section{Limitations and Future Directions}

\label{sec:limitations}

\subsection{Current Limitations}

\label{ssec:current_limits}

This work has several limitations worth acknowledging, and their impact on the reliability of the results is assessed for each:

\textbf{(i) Fixed design region size:} The $4.5\,\mu{\rm m}$ design side length is adequate for low topological charges ($\ell \le 2$), but for $\ell = 4, 5$ the theoretical LG spot peak radius $r_{\rm peak} = w_0\sqrt{\ell/2}$ is approaching or exceeding the design domain boundary, and the optimizer cannot fully express the radiation pattern of high-order modes within the finite space (see Section~\ref{ssec:high_order_limit} for details). Based on estimates, expanding the design domain to $6.5\,\mu{\rm m}$ or more could restore the mode purity for $\ell = 4, 5$ to above $85\%$.

\textbf{(ii) Single-wavelength design:} This work optimizes at a single target wavelength $\lambda = 0.593\,\mu{\rm m}$ (near the WS$_2$ exciton peak wavelength). GaP has relatively small dispersion near this wavelength ($\Delta n/\Delta\lambda \approx 0.3\,\mu{\rm m}^{-1}$), and purity is expected to remain above $80\%$ within about $\pm 5\%$ of the target wavelength, but this requires further simulation verification. For practical applications involving multi-wavelength OAM multiplexing, the FOM definition would need to be extended with multi-band constraints, or a broadband optimization strategy would need to be adopted.

\textbf{(iii) Ideal dipole source approximation:} The numerical model simulates WS$_2$ exciton radiation as an ideal circularly polarized point dipole ($100\%$ circular polarization purity), without accounting for three practical factors: (a) non-monochromaticity from exciton linewidth ($\sim$ tens of meV); (b) depolarization from inter-valley scattering at room temperature (experimental room-temperature valley polarization degree is typically $20\text{--}50\%$); (c) multi-exciton superposition effects. If the actual valley polarization degree $P_c$ is used as input, the far-field LG mode purity will approximately scale as ${\rm Purity}_{\rm exp} \approx P_c^2 \times {\rm Purity}_{\rm sim}$. This means that even with ${\rm Purity}_{\rm sim} = 98.66\%$, if $P_c = 50\%$, the experimentally expected purity would drop to about $25\%$; if Purcell enhancement raises $P_c \to 80\%$, it could recover to $\sim 63\%$. Improving the excitonic valley polarization purity is critical for actual device performance.

\textbf{(iv) Binary material assumption and fabrication tolerance:} The final structure is an ideally binarized GaP/air distribution. In actual fabrication, sidewall roughness ($\sim 5\,{\rm nm}$), etch depth errors ($\pm 5\%$), and GaP film refractive index non-uniformity could all introduce additional phase noise. Fabrication tolerance analysis should be an important part of future work, to establish a reliable mapping from numerical prediction to experimental measurement. Improvement directions for all of the above limitations will be proposed in Chapter 6.